\begin{papercard}{mayring2000}
\cf{Stages} \stagetag{C} \stagetag{I} \stagetag{L}. The English-language overview of Mayring's \emph{qualitative content analysis}: a rule-guided, category-centred method with two central procedures, \emph{inductive category development} and \emph{deductive category application}, a coding agenda, and a relaxed reliability standard ($\kappa>.7$). It is the origin of the German-language tradition that the later LLM paper (Mayring 2025) revisits.

\cf{Origin and aim} Developed about 20 years earlier in a longitudinal study of the psychosocial consequences of unemployment: about 600 open-ended interviews and more than 20,000 transcript pages \pc{mayring2000}{1}{Conducting about 600 open-ended interviews we received more than 20.000 pages of transcripts}. Aim: keep the strengths of quantitative content analysis (communication science) and carry them into qualitative-interpretative steps \pc{mayring2000}{1}{The main idea of the procedure of analysis is thereby, to preserve the advantages of quantitative content analysis}. Objects: any recorded communication (interviews, protocols, video, documents). Both manifest and latent content (Becker \& Lissmann) and formal aspects are in scope. It defines itself as ``empirical, methodological controlled analysis of texts within their context of communication, following content analytical rules and step by step models, without rash quantification'' \pc{mayring2000}{2}{Qualitative content analysis defines itself within this framework as an approach of empirical, methodological controlled analysis of texts}.

\cf{Four basic principles} (1) Fit the material into a \emph{model of communication}, deciding what the inferences are about (communicator, situation of production, socio-cultural background, the text, or effect). (2) \emph{Rules of analysis}: step by step, with content-analytical units. (3) \emph{Categories in the centre}, founded and revised through feedback loops. (4) \emph{Reliability and validity}: intersubjective comprehensibility, triangulation, and intercoder checks using trained project members, with a lowered threshold of Cohen's $\kappa>.7$ \pc{mayring2000}{3}{we use in qualitative content analysis (in contrary to quantitative content analysis) only trained members of the project team and we reduce the standard of coder agreement}.

\cf{Inductive category development} Categories are derived as close to the material as possible. Define a selection \emph{criterion} from theory and research question, work through the material, derive tentative categories step by step, revise in a feedback loop, reduce to main categories and check reliability (Figure 1). Quantitative analysis, such as category frequencies, is possible where the question suggests it \pc{mayring2000}{4}{The main idea of the procedure is, to formulate a criterion of definition, derived from theoretical background and research question}.

\cf{Deductive category application} Theory-derived categories are assigned to text passages in a methodologically controlled way (Figure 2). Each category gets \emph{explicit definitions, anchor examples and coding rules}, specifying exactly under which circumstances a passage is coded, and these are assembled into a \emph{coding agenda} (Figure 3: the self-confidence categories C1--C3, e.g.\ ``all three aspects of the definition have to point to high'', otherwise C2) \pc{mayring2000}{5}{Then main idea here is to give explicit definitions, examples and coding rules for each deductive category}. Agenda entries are completed step by step and revised during the analysis \pc{mayring2000}{6}{Category definitions, prototypical text passages, and rules for distinguishing different categories were formulated}.

\cf{Software and examples} Computers support, not replace, interpretation: assistant, documentation centre (to trace causes of non-reliability between coders), and a link to quantitative analysis \pc{mayring2000}{6}{The computer plays here a triple role}. Software named: ATLAS/ti, winMAX. Example studies: advisor theories (14 interviews, summarising), teacher diaries, media products, adult learners, Alzheimer biographies, unemployed teachers (50) combining inductive and deductive steps.

\cf{Limits} Less appropriate if the question is highly open-ended and categories would restrict, or if a holistic, non-step-wise analysis is intended. Combine with other procedures; research question and material come first \pc{mayring2000}{8}{The procedures of qualitative content analysis seem less appropriate}.

\cf{Use for this thesis} Supplies the German-tradition name for what the thesis does: deductive category application with a \emph{coding agenda} (category, definition, anchor example, coding rule) and a revision loop. The coding-agenda columns map one-to-one onto the codebook fields, and the ``borderline rule'' (e.g.\ when mixed evidence becomes the middle category) is a useful pattern for ordinal codes. The $\kappa>.7$ line is a published, citable lower threshold for exploratory work, to be justified against the stricter views (Hruschka; Krippendorff). It is a short overview, so details of the step models (Figures 1--2) must be taken from Mayring's German handbook if the thesis relies on them. Compare with the 2025 LLM paper before deciding whether LLM-assisted category application is in scope.
\end{papercard}
